\documentclass[12pt,a4paper]{article}
\usepackage[UTF8]{ctex}
\usepackage{geometry}
\usepackage{amsmath,amssymb}
\usepackage{array}
\usepackage{enumitem}
\usepackage{fancyhdr}
\usepackage{setspace}
\usepackage{longtable}
\usepackage{hyperref}

\geometry{left=2.5cm,right=2.5cm,top=2.5cm,bottom=2.5cm}
\setlength{\headheight}{15pt}
\setstretch{1.5}
\raggedbottom
\setlength{\emergencystretch}{3em}
\hfuzz=100pt
\hbadness=10000

\pagestyle{fancy}
\fancyhf{}
\fancyhead[C]{说明书}
\fancyfoot[C]{\thepage}

\newcommand{\Kearly}{K_{\mathrm{early}}}
\newcommand{\Kfinal}{K_{\mathrm{final}}}
\newcommand{\Kbase}{K_{\mathrm{base}}}
\newcommand{\Krein}{K_{\mathrm{reinforce}}}

\begin{document}

\begin{center}
    {\Huge\heiti 说明书}
\end{center}

\vspace{1em}

\noindent\textbf{发明名称：}一种面向6G网络切片的超低时延抗量子自适应握手方法、系统、设备及存储介质

\section*{一、技术领域}

本发明涉及移动通信安全、后量子密码学、网络切片安全编排以及低时延连接建立技术领域，尤其涉及一种面向 6G 网络切片场景的切片感知、能力感知、时延预算感知的抗量子自适应握手方法、系统、设备及存储介质。

\section*{二、背景技术}

随着量子计算技术的发展，传统基于 RSA、ECC 的公钥密码体制面临被量子算法破解的风险。后量子密码学（Post-Quantum Cryptography, PQC）为移动通信系统提供了量子韧性基础，但现有 PQC 算法通常具有公钥大、密文长、计算量高的特点，直接用于移动通信接入握手时会显著增加连接建立时延、空口占用和终端算力负担。

在 6G 网络中，不同网络切片面向的业务类型存在明显差异。例如，自动驾驶、工业控制和车路协同场景对时延极其敏感；高清视频和沉浸式业务对吞吐和长期安全更敏感；大规模传感业务又强调功耗与轻量接入。因此，若对所有切片统一采用同一种后量子握手流程，则要么无法满足低时延切片的连接建立需求，要么无法兼顾高安全级别切片的长期安全性。

现有技术中，一类方案将 PQC 直接接入 TLS、DTLS、QUIC 或 IPsec 等协议栈中，使握手阶段直接交换完整的后量子密钥材料；另一类方案采用双 KEM 或多 KEM 的静态组合方式，以提升鲁棒性；还有一类方案研究密码敏捷或动态算法选择机制。然而，这些方案通常存在以下问题：

\begin{enumerate}[leftmargin=2em]
    \item 对不同业务切片缺乏差异化的握手计划，无法根据切片 SLA 动态调整握手开销；
    \item 仍然采用单次握手中一次性交换大量后量子材料的方式，导致首包时延偏高；
    \item 即使采用多 KEM 或混合 KEM，也多为静态组合，未能将后量子安全增强过程拆分为多阶段执行；
    \item 切片安全、切片特定认证和实际后量子握手状态机之间缺少一体化设计；
    \item 对切片标识、握手计划、阈值参数和业务权限范围缺少统一的 transcript 绑定，易受降级攻击或跨切片重放攻击影响。
\end{enumerate}

因此，亟需一种新的抗量子握手技术，使得 6G 网络能够在满足不同切片安全要求的同时，降低连接建立时延并提升工程可部署性。

\section*{三、发明内容}

\subsection*{（一）发明目的}

本发明的目的在于提供一种面向 6G 网络切片的超低时延抗量子自适应握手方法、系统、设备及存储介质，以解决现有后量子握手无法兼顾切片差异化、超低时延和抗降级安全的问题。

\subsection*{（二）技术方案}

为实现上述目的，本发明提供如下技术方案。

一种面向 6G 网络切片的超低时延抗量子自适应握手方法，包括：

\begin{enumerate}[leftmargin=2em]
    \item 接收终端发送的握手请求消息，所述握手请求消息至少携带切片标识、终端能力信息以及终端支持的后量子密码算法集合；
    \item 根据所述切片标识对应的切片握手画像，并结合所述终端能力信息和当前链路状态信息，选择与目标切片对应的握手计划；
    \item 按照所述握手计划执行基础后量子密钥协商，获得基础份额，并导出基础会话密钥；
    \item 在允许早期业务的情况下，基于所述基础份额导出早期业务密钥；
    \item 生成增强密钥材料，并按照分片参数将所述增强密钥材料编码为多个增强片段；
    \item 在增强阶段时限内传输所述多个增强片段，并在接收到达到阈值的有效增强片段时恢复所述增强密钥材料；
    \item 基于所述基础份额和所述增强密钥材料导出最终会话密钥；
    \item 将切片标识、握手计划标识、算法包摘要、分片参数和增强阶段时限参数绑定到握手转录文本中。
\end{enumerate}

在一个实施方式中，所述切片握手画像至少包括：切片时延预算、最小安全等级、允许的算法集合、是否允许早期业务、增强片段恢复阈值、增强阶段失效窗口和业务权限级别中的至少一项。

在一个实施方式中，所述握手计划通过目标函数进行选择：
\[
b^\star=\arg\min_{b \in \mathcal{B}} \left(\alpha \hat{T}(b)+\beta \hat{C}(b)+\gamma \hat{B}(b)+\delta \hat{P}_{loss}(b)\right)
\]
其中，$\hat{T}(b)$ 表示预测握手完成时间，$\hat{C}(b)$ 表示计算成本，$\hat{B}(b)$ 表示握手字节开销，$\hat{P}_{loss}(b)$ 表示因链路丢包导致增强阶段超时的风险。

在一个实施方式中，所述基础后量子密钥协商可包括单一后量子 KEM、多个后量子 KEM 的组合、椭圆曲线密钥交换与后量子 KEM 的混合方式中的至少一种。

在一个实施方式中，所述增强密钥材料在分片前执行全有或全无变换、秘密共享、纠删编码中的至少一种，以降低部分分片泄露信息的风险。

在一个实施方式中，所述多个增强片段被映射到后续握手报文、首批业务报文、独立控制报文、多路径子流中的至少一种进行承载。

在一个实施方式中，所述早期业务密钥仅用于承载白名单业务，且白名单业务范围由切片策略限定，并被纳入握手转录文本绑定内容。

在一个实施方式中，若在增强阶段时限内未达到恢复阈值，则根据切片策略执行限制业务范围、降低会话权限、发起补强重协商、请求重发增强片段、撤销当前会话中的至少一种操作。

本发明还提供一种面向 6G 网络切片的超低时延抗量子自适应握手系统，包括：

\begin{enumerate}[label=\arabic*), leftmargin=2em]
    \item 切片策略控制模块，用于存储不同切片对应的切片握手画像；
    \item 握手计划决策模块，用于根据切片标识、终端能力信息和链路状态信息选择握手计划；
    \item 基础份额协商模块，用于执行基础后量子密钥协商并导出基础会话密钥；
    \item 早期业务控制模块，用于导出早期业务密钥并限制其业务承载范围；
    \item 增强材料编码模块，用于生成增强密钥材料并将其编码为满足阈值恢复条件的多个增强片段；
    \item 增强片段调度模块，用于调度增强片段的发送、接收、重传或多路径分发；
    \item 增强材料恢复模块，用于在接收到达到阈值的有效增强片段时恢复增强密钥材料；
    \item 最终密钥导出模块，用于导出最终会话密钥；
    \item 转录绑定保护模块，用于将关键协商参数绑定到握手转录文本中。
\end{enumerate}

本发明还提供一种电子设备以及一种计算机可读存储介质，用于执行上述方法。

\subsection*{（三）有益效果}

与现有技术相比，本发明至少具有如下有益效果：

\begin{enumerate}[leftmargin=2em]
    \item 能够根据不同网络切片的时延预算和安全需求动态选择后量子握手计划，避免所有业务承担统一且过高的握手开销；
    \item 通过“基础份额 + 增强份额”的分阶段设计，使连接首阶段更快建立可用安全通道；
    \item 通过对增强密钥材料进行阈值分片恢复，提高在无线链路丢包和重排序环境下的成功率；
    \item 通过对白名单早期业务和高敏业务进行差异化承载，提高协议在 URLLC、eMBB、mMTC 等多种切片中的适用性；
    \item 通过将切片标识、握手计划、分片参数等内容绑定到握手转录文本中，提高抗降级和抗跨切片重放能力；
    \item 适合部署于终端、边缘安全锚点、切片安全网关、核心网安全功能或接入代理中，具备较好的工程落地性和标准化潜力。
\end{enumerate}

\section*{四、附图说明}

为了更清楚地说明本发明实施例或现有技术中的技术方案，下面对实施例中所需要使用的附图作简单说明。所述附图仅为示意性附图。

\begin{enumerate}[leftmargin=2em]
    \item 图1为本发明面向 6G 网络切片的超低时延抗量子自适应握手系统架构示意图；
    \item 图2为本发明中切片握手画像选择和握手计划生成流程示意图；
    \item 图3为本发明中基础份额、增强份额、早期业务密钥和最终会话密钥之间的时序关系示意图；
    \item 图4为本发明中增强密钥材料的阈值分片与恢复流程示意图；
    \item 图5为本发明中自动驾驶切片、视频流切片和传感器切片的参数差异示意图。
\end{enumerate}

\section*{五、具体实施方式}

下面结合附图和具体实施方式对本发明作进一步详细说明。应当理解，这里描述的实施方式仅用于解释本发明，并不用于限定本发明。

\subsection*{（一）术语定义}

在本说明书中：

\begin{itemize}[leftmargin=2em]
    \item ``切片标识''记为 $SID$，用于标识所访问的网络切片或切片实例；
    \item ``切片握手画像''记为 $\Pi(SID)$，用于描述切片的安全和时延约束；
    \item ``基础份额''记为 $z_0$，指基础后量子协商阶段得到的共享秘密；
    \item ``增强密钥材料''记为 $\Krein$，指在基础阶段后继续补充的安全增强材料；
    \item ``基础会话密钥''记为 $\Kbase$；
    \item ``早期业务密钥''记为 $\Kearly$；
    \item ``最终会话密钥''记为 $\Kfinal$。
\end{itemize}

\subsection*{（二）系统架构实施方式}

如图1所示，本发明实施例中的系统至少包括终端设备 100、边缘安全锚点 200、切片策略控制器 300 和业务服务端 400。

终端设备 100 用于发起握手请求并提供终端能力信息。边缘安全锚点 200 或业务服务端 400 用于根据切片策略控制器 300 下发的切片握手画像，选择合适的握手计划并执行握手过程。

切片策略控制器 300 至少维护如下参数：

\begin{itemize}[leftmargin=2em]
    \item 切片最小时延预算；
    \item 切片最小安全等级；
    \item 切片允许的算法集合；
    \item 是否允许早期业务；
    \item 增强片段分片数与恢复阈值；
    \item 增强阶段时限；
    \item 异常处理策略。
\end{itemize}

\subsection*{（三）握手计划选择实施方式}

如图2所示，在一个实施方式中，终端首先发送握手请求消息，请求消息中包含：

\begin{itemize}[leftmargin=2em]
    \item 切片标识；
    \item 终端能力等级；
    \item 支持的 KEM 或混合密钥交换集合；
    \item 目标业务等级；
    \item 可选的切片恢复票据。
\end{itemize}

边缘安全锚点或服务端接收到请求后，读取切片握手画像 $\Pi(SID)$，并结合当前链路状态信息 $Net$ 以及终端能力向量 $Cap_{UE}$，在候选握手计划集合中选择最优计划。

在一个实施方式中，对于自动驾驶切片，系统更倾向于选择：

\begin{itemize}[leftmargin=2em]
    \item 基础阶段负担较轻的后量子协商方式；
    \item 较少的增强片段；
    \item 较短的增强阶段时限；
    \item 较严格的白名单业务范围。
\end{itemize}

对于高清视频切片，则可选择：

\begin{itemize}[leftmargin=2em]
    \item 更高安全级别的基础算法；
    \item 更多的增强片段；
    \item 更宽的增强阶段窗口；
    \item 不允许或较少允许早期业务。
\end{itemize}

\subsection*{（四）分阶段密钥形成实施方式}

如图3所示，本发明并不要求在握手首阶段一次性完成全部高开销后量子材料交换，而是将会话密钥形成拆分为基础阶段和增强阶段。

在基础阶段，系统执行基础后量子密钥协商，获得基础份额 $z_0$。然后按照如下方式导出基础会话密钥：
\[
\Kbase=\mathrm{HKDF}(z_0, H(\mathrm{transcript}), SID, \Pi(SID)).
\]

若该切片允许有限早期业务，则进一步导出：
\[
\Kearly=\mathrm{HKDF}(z_0, H(\mathrm{transcript}), SID, \texttt{plan\_id}, \texttt{early\_scope}).
\]

在增强阶段，系统生成一个或多个增强份额 $\{z_1,\dots,z_m\}$，并计算：
\[
\Krein = H(z_1 \,\|\, z_2 \,\|\, \cdots \,\|\, z_m).
\]

之后，再导出最终会话密钥：
\[
\Kfinal=\mathrm{HKDF}(z_0 \,\|\, \Krein \,\|\, H(\mathrm{transcript}), SID, \texttt{plan\_id}).
\]

\subsection*{（五）增强材料阈值分片实施方式}

如图4所示，为避免增强阶段成为新的单点等待瓶颈，本发明将增强密钥材料 $\Krein$ 编码为多个增强片段 $r_1,\dots,r_n$，满足任意达到阈值 $t$ 的有效增强片段即可恢复所述增强密钥材料。

在一个实施方式中，采用如下处理流程：

\begin{enumerate}[leftmargin=2em]
    \item 对增强密钥材料执行全有或全无变换；
    \item 对变换结果执行秘密共享或纠删编码；
    \item 生成总数为 $n$ 的增强片段；
    \item 将各增强片段分配到后续握手报文、控制报文、首批数据报文或多路径子流中。
\end{enumerate}

在无线环境中，即使存在片段丢失或重排序，只要在增强阶段时限内收到至少 $t$ 个有效增强片段，即可恢复增强密钥材料并切换至最终会话密钥。

\subsection*{（六）业务分级承载实施方式}

在一个实施方式中，本发明并不允许所有业务在基础阶段完成后立即全部放行，而是采用业务分级承载机制。

对于自动驾驶切片，只允许少量启动控制消息在 $\Kearly$ 上承载；对于工业控制切片，允许状态登记类消息先行，而正式控制命令必须等待 $\Kfinal$ 建立后才允许发送；对于高清视频业务，可以直接禁止早期业务，等待最终会话密钥建立后再开始大流量传输。

所述早期业务范围由切片策略预定义，并通过 \texttt{early\_scope} 字段写入握手转录文本绑定内容中。

\subsection*{（七）抗降级与跨切片隔离实施方式}

为防止攻击者通过修改协商内容诱导双方退回到更弱但更快的握手模式，本发明要求将以下内容绑定进握手转录文本：

\begin{itemize}[leftmargin=2em]
    \item 切片标识；
    \item 握手画像标识；
    \item 握手计划标识；
    \item 算法包摘要；
    \item 增强片段总数；
    \item 恢复阈值；
    \item 增强阶段时限；
    \item 早期业务范围。
\end{itemize}

通过上述绑定，可实现：

\begin{enumerate}[leftmargin=2em]
    \item 抗降级攻击；
    \item 抗跨切片重放攻击；
    \item 在多路径增强片段传输场景下的路径篡改检测。
\end{enumerate}

\subsection*{（八）切片恢复票据实施方式}

在一个实施方式中，系统可生成与切片绑定的恢复票据，用于后续快速重连。恢复票据可表示为：
\[
Ticket = \mathrm{AEAD}_{K_S}(SID, \texttt{qprofile\_id}, \texttt{plan\_id}, Cap_{UE}, Expiry, PolicyHash).
\]

在终端再次进入相同切片时，可以通过所述恢复票据快速恢复切片握手画像和握手计划，减少重复协商开销，但该票据不能跨切片复用。

\subsection*{（九）实施例1：自动驾驶切片}

在本实施例中，目标切片为超可靠低时延通信切片。系统为该切片选择较轻量的基础阶段算法和较少的增强片段，允许极少量启动控制消息在早期业务密钥上承载。一旦增强片段达到恢复阈值，即立即切换到最终会话密钥，以满足超低时延连接建立需求。

\subsection*{（十）实施例2：高清视频切片}

在本实施例中，目标切片为增强移动宽带切片。系统为该切片选择更高安全等级的算法集合和较多的增强片段，并允许更长的增强阶段时限。由于该切片对首包时延不如 URLLC 场景敏感，因此可以优先保证长会话安全性和抗量子冗余。

\subsection*{（十一）实施例3：工业传感切片}

在本实施例中，目标切片为大规模机器类通信切片。终端设备算力受限，因此系统根据终端能力等级降低基础阶段计算负担，并由边缘安全锚点承担更多增强片段管理逻辑。正式控制命令仅在最终会话密钥建立后允许下发。

\subsection*{（十二）电子设备实施方式}

本发明还提供一种电子设备，包括处理器和存储器，所述存储器中存储有计算机程序，所述处理器执行所述计算机程序时，用于实现上述任一实施方式中的方法。

\subsection*{（十三）存储介质实施方式}

本发明还提供一种计算机可读存储介质，其上存储有计算机程序，所述计算机程序被处理器执行时，用于实现上述任一实施方式中的方法。

\subsection*{（十四）等同替换说明}

本领域技术人员在不脱离本发明精神和保护范围的情况下，对本发明所作出的任何等同替换、改进和变形，均应落入本发明的保护范围内。

\end{document}
